% Charges partielles du bromure d'hydrogène et de la méthylamine (corrigé de l'exercice 2, chapitre 5)
{\setchemfig{atom sep=2.1em}\renewcommand{\arraystretch}{1}%
\begin{tabular}{@{}c@{\hspace{18mm}}c@{}}
\chemfig{@{xa}H-@{xb}\charge{90=\|,0=\|,270=\|}{Br}}%
\chemmove{\node[pcDelta,anchor=south] at ([yshift=1pt]xa.north) {\dplus};
  \node[pcDelta,anchor=south] at ([yshift=4.5pt]xb.north) {\dmoins};}
&
\chemfig{@{ya}C(-[2]H)(-[4]H)(-[6]H)-@{yb}\charge{90=\|}{N}(-[6]@{yc}H)-@{yd}H}%
\chemmove{\node[pcDelta,anchor=south west] at ([xshift=1pt,yshift=1pt]ya.north east) {\dplus};
  \node[pcDelta,anchor=south] at ([yshift=4.5pt]yb.north) {\dmoins};
  \node[pcDelta,anchor=west] at ([xshift=1pt]yc.east) {\dplus};
  \node[pcDelta,anchor=south] at ([yshift=1pt]yd.north) {\dplus};}\\[1mm]
{\otLight\fontsize{7.5}{9}\selectfont\color{otMuted}bromure d'hydrogène} &
{\otLight\fontsize{7.5}{9}\selectfont\color{otMuted}méthylamine}
\end{tabular}}
